L'acide nitrique \ce{HNO3} est un liquide moléculaire à caractère dipolaire~; sa
dissolution dans l'eau libère des ions oxonium et des ions nitrate.
\begin{enumerate}
    \item Écrire l'équation de la dissolution de l'acide nitrique dans l'eau.
    \item Sur un flacon d'acide nitrique, on relève les informations suivantes~:
          \begin{center}
              acide nitrique~:  68\%~; $d=1,41$.
          \end{center}

          Calculer la concentration molaire de la solution contenue dans le
          flacon.
    \item Quel volume de cette solution doit-on prélever pour préparer
          $\qty{1.0}{\L}$ de solution telle que~:
          $[\ce{NO$_{3}^{-}$(aq)}]=\qty{0.15}{\mol\per\L}$~?

\end{enumerate}

